% Generated draft: requires mathematical and source review.
\documentclass[11pt]{article}
\usepackage[margin=25mm]{geometry}
\usepackage{fontspec}
\setmainfont{Latin Modern Roman}
\usepackage{amsmath,amssymb,mathtools,unicode-math}
\setmathfont{Latin Modern Math}
\usepackage{enumitem,needspace,xurl,hyperref}
\setlength{\parindent}{0pt}
\setlength{\parskip}{5pt}
\setlength{\emergencystretch}{3em}
\newcommand{\N}{\mathbb{N}}
\newcommand{\Z}{\mathbb{Z}}
\newcommand{\Q}{\mathbb{Q}}
\newcommand{\R}{\mathbb{R}}
\newcommand{\C}{\mathbb{C}}
\newcommand{\F}{\mathbb{F}}
\newcommand{\A}{\mathbb{A}}
\newcommand{\PP}{\mathbb{P}}
\newcommand{\E}{\mathbb{E}}
\newcommand{\eps}{\varepsilon}
\newcommand{\dd}{\,\mathrm{d}}
\begin{document}
\section*{UnsolvedMath ID: 10600013}
\textbf{Source problem code:} catalogue entry
\subsection{Definitions and mathematical statement}
For a virtual knot diagram \(D\), let \(c_{\mathrm{cl}}(D)\) be its number of classical crossings, with virtual crossings excluded. Let \(B(D)\) denote its all-\(B\) smoothing state, subject to the source's unspecified smoothing convention. Write \(\langle D\rangle\) for the Kauffman bracket and \(\operatorname{Kh}(K;\Z)\) for the integral Khovanov theory intended by the source when defined.

\emph{Compound problem: Khovanov homology.}
\begin{enumerate}
\item Construct a generalization of the Khovanov complex that works for arbitrary virtual diagrams, and investigate the existing construction referred to in the source. That construction uses an orientable-atom condition to obtain integral homology for large classes of virtual links. Seek a theory over \(\Z\) for all virtual knots, with diagram invariance under the virtual moves. The source also asks analogous questions for evolving homology theories related to other quantum invariants. Its citations and the exact existing complex are missing.
\item Classify all \(K\)-full virtual knots and all \(\mathrm{Kho}\)-full virtual knots. The source calls a knot \(K\)-full when some diagram has its leading term, its lowest term, or both coming from the \(B\)-state, and defines \(\mathrm{Kho}\)-full analogously relative to the Khovanov invariant. It calls the witnessing diagrams optimal and observes that examples failing both fullness conditions are easy to find. The three alternatives ``leading,'' ``lowest,'' and ``both'' are not identified as one precise condition in the excerpt; the relevant version and meaning of ``comes from'' must be fixed before a classification is well-defined.
\item Are optimal diagrams always minimal in classical crossing number among virtual diagrams of the same virtual knot? With \(c_v(K)=\min_{D\sim K}c_{\mathrm{cl}}(D)\), the question is whether every intended optimal \(D\) satisfies \(c_{\mathrm{cl}}(D)=c_v(K)\).
\item Classify all diagram moves preserving optimality. The allowed class of moves, and whether preservation is required universally or in specified configurations, are not stated.
\item If a classical knot \(K\) has minimal classical crossing number \(n\), must every virtual diagram \(D\) representing \(K\) satisfy \(c_{\mathrm{cl}}(D)\geq n\)? Equivalently, is \(c_v(K)=n\) for every classical \(K\), with the classical-to-virtual identification intended by the source? The minimum defining \(n\) is over classical diagrams; the universal quantifier here is over virtual diagrams.
\item Can a virtual knot have torsion in the \(B\)-state of its Khovanov homology, described in the source as the genuine leading term of some diagram? In the ordinary existential reading, ask whether there are a knot, diagram, and a nonzero finite-order class in the intended \(B\)-state-associated homology group. The source expressly specifies its cited formulation of Khovanov homology, but that citation and the relevant group or grading are absent. A chain-level state is not automatically a canonical summand of homology, so no such summand is presumed.
\end{enumerate}
The source says that the lowest and leading Kauffman-bracket terms for virtual knots were studied in several references, all lost from the supplied extraction. Its orientable-atom limitation is historical context, not a claim that no integral theory now exists.
\subsection{Short English statement}
The item combines integral Khovanov constructions for arbitrary virtual diagrams with classification of fullness and optimal diagrams, crossing-minimality questions, optimality-preserving moves, and possible torsion in a specified extreme \(B\)-state contribution. Its exact fullness conventions and several essential citations are missing.
\subsection{Sources}
\begin{itemize}
\item Local UnsolvedMath registry, permanent ID 10600013, source code catalogue entry, ``Virtual-knot problem 13 --- Khovanov Homology.'' No individual proposer is supplied.
\item \emph{Unsolved Problems in Virtual Knot Theory and Combinatorial Knot Theory}, revised survey, problem-list item 13. Supplied URL: \url{https://arxiv.org/abs/1409.2823}; recorded access 2026-07-29; extraction source-tex. Its dedicated section contains 60 top-level enumerated items and retains compound items intact.
\item The extracted statement has missing citations for the orientable-atom construction, related quantum-invariant homologies, the formulation used for \(B\)-state torsion, and studies of extreme bracket terms. No bibliographic entries can be recovered from this record.
\item Supplied status partially\_solved and PARTIAL-PROGRESS. The summary says integral virtual Khovanov homology exists through several constructions, naming Manturov, ``Khovanov--Rozansky over integers,'' and ``KLS for thickened surfaces,'' without full citations; it says the classification, optimality, and torsion questions appear open. These are supplied assessments. Current primary-literature status and redistribution terms for the publicly accessible author/publisher copy remain NEEDS\_REVIEW.
\end{itemize}
\end{document}
